% Copyright (c) 2025 Intel Corporation.  All rights reserved.  See the file COPYRIGHT for more information.
% SPDX-License-Identifier: Apache-2.0

\documentclass[11pt]{article}

% Page layout
\usepackage[letterpaper, margin=1in]{geometry}
\setlength{\headheight}{14pt}
\usepackage{parskip}
\setlength{\parindent}{0pt}

% Typography
\usepackage{fancyvrb}
\usepackage{booktabs}

% Code listings
\usepackage{listings}
\lstset{
    basicstyle=\ttfamily\small,
    columns=fullflexible,
    keepspaces=true,
    showstringspaces=false,
    breaklines=true,
    frame=none,
    xleftmargin=2em,
}

% For BNF-style syntax rules
\usepackage{syntax}
\setlength{\grammarparsep}{4pt}
\setlength{\grammarindent}{11em}

% Hyperlinks
\usepackage[colorlinks=true, linkcolor=blue, urlcolor=blue, citecolor=blue]{hyperref}

% Headers and footers
\usepackage{fancyhdr}
\pagestyle{fancy}
\fancyhf{}
\fancyhead[L]{\textit{The cspbuild System Builder}}
\fancyhead[R]{\thepage}
\renewcommand{\headrulewidth}{0.4pt}

% Custom commands (matching csp.tex)
\newcommand{\keyword}[1]{\texttt{\textbf{#1}}}
\newcommand{\nonterminal}[1]{\textit{#1}}

\title{\textbf{The cspbuild System Builder} \\[0.5em]
\large User Manual for the \texttt{.sys} Format}
\author{Mika Nystr\"om \\
  Neuromorphic Computing Laboratory \\
  Intel Corporation}
\date{DRAFT Version 0.1\\
\today}

\begin{document}

\maketitle

\vspace{1em}

\noindent
 Copyright \copyright\ 2025 Intel Corporation

 Licensed under the Apache License, Version 2.0 (the ``License'');
 you may not use this file except in compliance with the License.
 You may obtain a copy of the License at

 {\tt http://www.apache.org/licenses/LICENSE-2.0}

\section{Introduction}

The \texttt{cspbuild} system builder automates the construction of
multi-process CSP simulations.  It replaces the manual workflow of
running \texttt{cspfe}, writing \texttt{.procs} files, patching
cellinfo, and invoking \texttt{drive!}\ with a single command.

The input is a \texttt{.sys} file that declares the channels, process
types, and instances of a system.  Each process type refers to a
\texttt{.csp} source file (or contains inline CSP code).  The builder
parses and validates the system description, runs \texttt{cspfe} on
each process, patches the resulting S-expression files with channel
port definitions, writes the \texttt{.procs} file, and calls
\texttt{drive!}\ and \texttt{cm3} to produce a simulator binary.

\section{Quick Start}

\subsection{A single-process system}

Create a file \texttt{hello.csp}:
\begin{Verbatim}[samepage=true]
print("hello, world!")
\end{Verbatim}

Create a file \texttt{hello.sys}:
\begin{Verbatim}[samepage=true]
system Hello;
  process HELLO = "hello.csp";
begin
  var x : HELLO;
end.
\end{Verbatim}

Build and run:
\begin{Verbatim}[samepage=true]
$ cspc -scm /dev/stdin <<'EOF'
(load "$M3UTILS/csp/src/setup.scm")
(load "$M3UTILS/csp/src/cspbuild.scm")
(build-system! "hello.sys")
(exit 0)
EOF
$ build/ARM64_DARWIN/sim
1: x.x: hello, world!
\end{Verbatim}

\subsection{A producer-consumer system}

Create \texttt{producer.csp}:
\begin{Verbatim}[samepage=true]
int x;
*[ R!x ; x = x + 1 ]
\end{Verbatim}

Create \texttt{consumer.csp}:
\begin{Verbatim}[samepage=true]
int x;
*[ L?x ; print(x) ]
\end{Verbatim}

Create \texttt{prodcons.sys}:
\begin{Verbatim}[samepage=true]
system ProdCons;
  var c : channel(32);

  process Producer = "producer.csp"
    port out R : channel(32);

  process Consumer = "consumer.csp"
    port in L : channel(32);
begin
  var p : Producer(R => c);
  var q : Consumer(L => c);
end.
\end{Verbatim}

Build and run:
\begin{Verbatim}[samepage=true]
$ cspc -scm /dev/stdin <<'EOF'
(load "$M3UTILS/csp/src/setup.scm")
(load "$M3UTILS/csp/src/cspbuild.scm")
(build-system! "prodcons.sys")
(exit 0)
EOF
$ build/ARM64_DARWIN/sim
3: x.q: 0x0
4: x.q: 0x1
5: x.q: 0x2
...
\end{Verbatim}

\section{The \texttt{.sys} File Format}

A \texttt{.sys} file describes a system of communicating CSP processes.
It declares the channels that connect them, the process types (each
backed by a \texttt{.csp} source file or an inline body), and the
instances that make up the system.

\subsection{Overall structure}

\begin{Verbatim}[samepage=true]
system <name>;
  <declarations>
begin
  <instances>
end.
\end{Verbatim}

The system name is used as a prefix for process type names in the
generated files.  For a system named \texttt{Hello} with a process
\texttt{HELLO}, the cell name becomes \texttt{Hello.HELLO} and the
escaped filename becomes \texttt{Hello\_46\_HELLO.scm}.

\subsection{Comments}

Two comment styles are supported:
\begin{itemize}
\item Line comments: \verb!// comment to end of line!
\item Block comments: \verb!(* comment *)!, which may be nested
\end{itemize}

\subsection{Declarations}

Declarations appear between \keyword{system} and \keyword{begin}.
There are two kinds: channel declarations and process declarations.

\subsubsection{Channel declarations}

\begin{Verbatim}[samepage=true]
var <name> : channel(<width>);
var <name> : channel(<width>) slack <depth>;
\end{Verbatim}

Declares a channel of the given bit width.  The optional
\keyword{slack} clause sets the channel buffer depth (default~1).

Examples:
\begin{Verbatim}[samepage=true]
var c : channel(32);            // 32-bit channel, slack 1
var wide : channel(64) slack 4; // 64-bit channel, slack 4
\end{Verbatim}

\subsubsection{Process declarations}

\begin{Verbatim}[samepage=true]
process <name> = <body> <port-declarations>;
\end{Verbatim}

The body is either a string literal naming an external \texttt{.csp} file:
\begin{Verbatim}[samepage=true]
process Producer = "producer.csp"
    port out R : channel(32);
\end{Verbatim}

or an inline CSP body enclosed in \keyword{begin}\ldots\keyword{end}:
\begin{Verbatim}[samepage=true]
process HELLO = begin print("hello!") end;
\end{Verbatim}

Inline bodies are written to a temporary file and passed to
\texttt{cspfe}.  They may contain nested
\keyword{begin}\ldots\keyword{end} blocks.

\subsubsection{Port declarations}

Port declarations follow the process body (before the terminating
semicolon) and declare the channels that the process communicates on.
Each port has a direction (\keyword{in} or \keyword{out}) and a
channel width that must match the channel it is bound to.

\begin{Verbatim}[samepage=true]
port in <name> : channel(<width>)
port out <name> : channel(<width>)
\end{Verbatim}

The port name must match the channel variable used inside the
\texttt{.csp} source.  For example, if the CSP code contains
\verb!L?x! and \verb!R!x!, the process declaration should have
\verb!port in L! and \verb!port out R!.

A process with no ports (such as a standalone ``hello world'' process)
simply has no port declarations.

\subsection{Instances}

Instances appear between \keyword{begin} and \keyword{end}.
Each instance creates one copy of a process type and binds its
ports to channels.

\begin{Verbatim}[samepage=true]
var <instance-name> : <process-name>;
var <instance-name> : <process-name>(<bindings>);
\end{Verbatim}

Bindings are comma-separated \nonterminal{port} \verb!=>!
\nonterminal{channel} pairs:

\begin{Verbatim}[samepage=true]
var p : Producer(R => c);
var q : Consumer(L => c);
\end{Verbatim}

The binding \verb!R => c! connects the process's port \texttt{R}
to the system-level channel~\texttt{c}.  Every port declared on
the process type must be bound in every instance.

\section{Formal Syntax}

\begin{grammar}
<system-file> ::= \keyword{system} \nonterminal{IDENT} ``;'' \{ <declaration> \}
\keyword{begin} <instance-list> \keyword{end} ``.''

<declaration> ::= <channel-decl> | <process-decl>

<channel-decl> ::= \keyword{var} \nonterminal{IDENT} ``:'' \keyword{channel}
``('' \nonterminal{INTEGER} ``)'' [ \keyword{slack} \nonterminal{INTEGER} ] ``;''

<process-decl> ::= \keyword{process} \nonterminal{IDENT} ``='' <process-body>
\{ <port-decl> \} ``;''

<process-body> ::= \nonterminal{STRING} | \keyword{begin} \nonterminal{csp-text} \keyword{end}

<port-decl> ::= \keyword{port} <direction> \nonterminal{IDENT} ``:'' \keyword{channel}
``('' \nonterminal{INTEGER} ``)''

<direction> ::= \keyword{in} | \keyword{out}

<instance-list> ::= \{ \keyword{var} \nonterminal{IDENT} ``:'' \nonterminal{IDENT}
[ ``('' <binding-list> ``)'' ] ``;'' \}

<binding-list> ::= <binding> \{ ``,'' <binding> \}

<binding> ::= \nonterminal{IDENT} ``\verb!=>!'' \nonterminal{IDENT}
\end{grammar}

\section{The \texttt{build-system!}\ Procedure}

\subsection{Synopsis}

\begin{Verbatim}[samepage=true]
(build-system! <sys-file> [option value] ...)
\end{Verbatim}

\subsection{Options}

\begin{tabular}{lll}
\toprule
Option & Default & Description \\
\midrule
\texttt{'slack} & \texttt{*default-slack*} & Channel buffer depth \\
\texttt{'run} & \texttt{\#f} & Run the simulator after building \\
\texttt{'verbose} & \texttt{\#f} & Print extra progress information \\
\bottomrule
\end{tabular}

\bigskip

Examples:
\begin{Verbatim}[samepage=true]
(build-system! "hello.sys")
(build-system! "prodcons.sys" 'slack 2 'verbose #t)
(build-system! "prodcons.sys" 'run #t)
\end{Verbatim}

\subsection{Other useful entry points}

The \texttt{cspbuild} library also exposes lower-level procedures for
use in interactive \texttt{cspc} sessions:

\begin{description}
\item[\texttt{(parse-sys-file \nonterminal{filename})}]
  Parse a \texttt{.sys} file and return its abstract syntax tree
  as a Scheme list.  Useful for debugging.

\item[\texttt{(validate-system! \nonterminal{sys})}]
  Check that all instances reference declared process types,
  all bindings reference declared ports and channels, all port
  widths match channel widths, and every port is bound.
  Raises an error on the first violation found.

\item[\texttt{(escape-name \nonterminal{string})}]
  Escape special characters in a name for use as a filename,
  matching the encoding used by \texttt{M3Ident.Escape}:
  \texttt{.}~$\to$~\texttt{\_46\_},
  \texttt{(}~$\to$~\texttt{\_40\_},
  \texttt{)}~$\to$~\texttt{\_41\_},
  \texttt{,}~$\to$~\texttt{\_44\_}.
\end{description}

\section{Build Pipeline}

The \texttt{build-system!}\ procedure executes the following steps:

\begin{enumerate}
\item \textbf{Parse} the \texttt{.sys} file into an abstract syntax tree.
\item \textbf{Validate} the system: check that all references are
  consistent and all ports are bound.
\item \textbf{Generate \texttt{.scm} files}: for each process type,
  run \texttt{cspfe -{}-scm -{}-name \nonterminal{cellname}
  \nonterminal{source.csp}} to produce an S-expression file.
  For processes with channel ports, read the file back, replace the
  cellinfo's port-definition list with one that includes the proper
  channel port definitions (sub-ports \texttt{C}, \texttt{q},
  \texttt{a}, \texttt{D}), and write it out again.
\item \textbf{Generate \texttt{.procs} file}: write one line per
  instance with its cell type, escaped name, and port-to-channel
  bindings.
\item \textbf{Compile}: call \texttt{drive!}\ (from
  \texttt{driver.scm}), which reads the \texttt{.procs} file,
  compiles all \texttt{.scm} files to Modula-3 via
  \texttt{compile-csp!}, and generates \texttt{SimMain.m3} and the
  \texttt{m3makefile}.
\item \textbf{Build}: run \texttt{cm3 -build -override} in the
  generated \texttt{build/src} directory to produce the simulator
  binary.
\item \textbf{Run} (optional): if the \texttt{'run} option is set,
  execute the simulator.
\end{enumerate}

\section{Limitations}

\begin{itemize}
\item All \texttt{.csp} files referenced by external process bodies
  must be accessible from the working directory where
  \texttt{build-system!}\ is invoked.
\item The \texttt{cspfe} binary must be built for the current platform.
  Its path is derived from the \texttt{M3UTILS} environment variable
  as \texttt{\$M3UTILS/csp/cspparse/ARM64\_DARWIN/cspfe}.
\item Channel types are currently limited to bundled-data channels
  (\texttt{standard.channel.bd}).
\item The \texttt{.sys} format does not support parameterized process
  types (as CAST does with e.g.\ \texttt{P0(int WIDTH)}).
  Each distinct parameterization requires a separate process
  declaration.
\end{itemize}

\end{document}
